\section{Use later failures to check what was actually published}
Prepare \texttt{s27-tight} from your completed collector and use the ordinary worked
file. The file loader still succeeds with 115 bytes. The collector's selected compressed
capacity is now 45. Its first IDAT can reserve and copy one byte; only 44 remain when
the second IDAT requests 45. The call fails with status one and returned count zero.
The destination may contain that first \texttt{78}; all four descriptor output fields
must still be sentinels. A partial scratch prefix is not a successful joined stream.

Next create \texttt{bad-crc.png} and use \texttt{s27-bad-crc}. Only the final stored
IEND CRC has been damaged. Earlier chunks can already have produced all 46 compressed
bytes, but the container still fails. Inspect the descriptor at
\texttt{collector\_return}: no width, height or length output was committed. This is
why the publication region belongs after IEND validation rather than inside IHDR.

Finally create \texttt{idat-gap.png} and use \texttt{s27-idat-gap}. This input inserts
a correctly checksummed tEXt between the two IDAT chunks. The metadata itself is
skippable, but its position ends the consecutive IDAT sequence. The later IDAT must
fail. Checking every CRC cannot replace checking chunk order.

In the supplied program these are PNG-stage failures, so it prints
\texttt{Invalid or unsupported PNG.} A successful collection instead reaches the
still-unfinished inflater boundary and prints the development message. Neither path
publishes framebuffer words or final decoder-context outputs. Inspect the collector's
own status at its return checkpoint to distinguish its success from that later
development status.
